\documentclass[12pt]{article}
\usepackage{amsfonts,amsthm,amsmath,amssymb,mathtools}
\usepackage{mathrsfs}
\usepackage{enumitem}
\usepackage{tikz-cd}
\usepackage{geometry}
\usepackage{mlmodern}

\geometry{letterpaper, portrait, margin=1in}

\setcounter{section}{0}

\renewcommand{\thesection}{\S\arabic{section}}
\renewcommand{\thesubsection}{\arabic{section}}
\newcommand{\dd}{\mathop{}\!\mathrm{d}}

\newtheorem{thm}{Theorem}
\newenvironment{theorem}[1]{
	\renewcommand{\thethm}{#1}
	\begin{thm}
		}{\end{thm}}

\newtheorem{lma}{Lemma}
\newenvironment{lemma}[1]{
	\renewcommand{\thelma}{#1}
	\begin{lma}
		}{\end{lma}}

\theoremstyle{definition}
\newtheorem{ex}{}
\newenvironment{exercise}[1]{
	\renewcommand{\theex}{#1}
	\begin{ex}
		}{\end{ex}}

\title{Homework 6}
\author{Connor Mooneyhan}
\date{October 10, 2025}

\begin{document}

\maketitle

\begin{ex}
	Let $E/\mathbb{C}$ be an elliptic curve and $P \in E(\mathbb{C})$ be a point.
	Prove that there are $m^{2}$ points $Q \in E(\mathbb{C})$ so that $mQ = P$.
\end{ex}
\begin{proof}
  We will work within the fundamental parallelogram, simplified further to act like a square for cleanliness.
  We recognize that any points ``found'' would need to be mapped back onto the curve to truly count as points on the curve.

  We identify each point $P$ with its corresponding point on the fundamental parallelogram $P_x\omega_1 + P_y\omega_2$ under an appropriate fixed isomorphism, referring to $P$ as just $(P_x, P_y)$ where $0 \leq x < 1$ and $0 \leq y < 1$.
  Under the hood, you could view this simplification as taking another isomorphism, this time an isomorphism from the vector space generated by $\omega_1$ and $\omega_2$ to the ``square'' vector space generated by $(0, 1)$ and $(1, 0)$.
  Thus, we interpret $mP$ as simply the linear multiple $(mP_x, mP_y)$.
  Framed this way, the question becomes much simpler.
  We are forced by the equation $mQ = P$ to have $mQ_x - P_x \in \mathbb{Z}$ and $mQ_y - P_y \in \mathbb{Z}$.
  This is easy enough!
  We just take $Q$ to be $(\frac{a + P_x}{m}, \frac{b + P_y}{m })$ where $a$ and $b$ are integers between 0 and $m-1$, inclusive.
  Since each coordinate can take on $m$ values, there are $m^2$ such points.
\end{proof}

\begin{ex}
	Let $E : y^{2} = x^{3} - 24003x + 1296702$.
	Use the Lutz-Nagell theorem to compute the set of points of finite order on $E$.
	You may use Magma to do computations, but you cannot just ask Magma for the torsion subgroup and have that suffice.
\end{ex}
\begin{proof}
	By Nagell-Lutz, we can reduce this to two cases, namely where $y = 0$ and where $y$ divides the discriminant $D$ (which Magma tells us is $158687432294400$).

	If $y = 0$, then we are just finding roots of the cubic on the right hand side of the equation.
	Magma tells us that these roots are approximately $-267.9$, $133.9 - 172.7i$, and $133.9 + 172.7i$, none of which are integers, and thus none of which are $x$-coordinates of rational points of finite order.

	Now we turn to divisors of the discriminant, which as mentioned earlier, Magma tells us is $158687432294400$, whose prime factorization is $2^{14} \cdot 3^{18} \cdot 5^2$ (again, according to Magma).
	Thus any integer $2^a \cdot 3^b \cdot 5^c$ is an option for $y$ where $0 \leq a \leq 14$, $0 \leq b \leq 18$, and $0 \leq c \leq 2$.
	I checked every single one of these with Magma using the code at the end of this answer, and I found no integers, which would lead me to conclude that there are no rational points of finite order.
	Of course, there's probably some mistake here, but I'm not sure what it is.

  \begin{verbatim}
	C<i> := ComplexField();
	PC<x> := PolynomialRing(C);

	for a in [0..14] do;
    for b in [0..18] do;
      for c in [0..2] do;
        y := 2^a * 3^b * 5^c;
        roots := Roots(x^3 - 24003 * x + 12796702 - y^2);
        root1 := roots[1][1];
        root2 := roots[2][1];
        root3 := roots[3][1];

        if IsIntegral(root1) then;
          root1;
        end if;
        if IsIntegral(root2) then;
          root2;
        end if;
        if IsIntegral(root3) then;
          root3;
        end if;
      end for;
    end for;
	end for;
  \end{verbatim}
\end{proof}

\begin{ex}
	$*$ Let $E : y^{2} = x^{3} + ax^{2} + bx + c$ be an elliptic curve with $a, b, c \in \mathbb{Z}$ and let $D$ be the discriminant.
	The duplication formula gives that if $(x,y)$ is a point on $E$, then the $x$-coordinate of $2P$ is \begin{equation*}
		\frac{\phi(x)}{4 f(x)} = \frac{x^{4} - 2bx^{2} - 8cx + b^{2} - 4ac}{4x^{3} + 4ax^{2} + 4bx + 4c}.
	\end{equation*}
	\begin{enumerate}[label=(\alph*)]

		\item Show that there are polynomials $F(x)$ and $\Phi(x)$ with integer coefficients so that \begin{equation*}
			      F(x) f(x) + \Phi(x) \phi(x) = D.
		      \end{equation*}
		      (Feel free to use Magma on this part.)

		\item Use (a) to show that if $P = (x,y)$ is a point with finite order, then
		      either $y = 0$ or $y^{2} | D.$
	\end{enumerate}
\end{ex}
\begin{proof}\
  \begin{enumerate}[label=(\alph*)]
    \item I'm really not sure how one would do this, even with Magma.
    \item We already have the $y = 0$ part from Nagell-Lutz, so we want to show that if $y | D$, then $y^2 | D$.
      I haven't quite figured out how to connect the dots here, but I have a feeling that it would be helpful to recognize that $f(x) = y^2$.
      This gets us to the equation \begin{equation*}
        F(x)y^2 + \Phi(x)\phi(x) = D.
      \end{equation*}
      I'm not sure how this gets us to $y^2$ dividing $D$ though, since it's not clear at all on the face of it if you could factor out a $y^2$ from $\Phi(x)\phi(x)$.
  \end{enumerate}
\end{proof}

\end{document}
